\documentclass[../DoAn.tex]{subfiles}
\begin{document}

Chương này trình bày các nền tảng lý thuyết và công nghệ được sử dụng để xây dựng hệ thống multi-tenant RAG chatbot hỗ trợ tư vấn bán hàng nội thất. Các nội dung được lựa chọn dựa trên những yêu cầu đã xác định ở Chương 2, bao gồm tư vấn sản phẩm theo cửa hàng, quản lý chatbot và knowledge base, quản lý hội thoại, lead, yêu cầu mua hàng, tham chiếu giá cơ bản, tích hợp kênh chat ngoài và quản lý vận hành hệ thống.

Mục tiêu của chương không phải là trình bày chi tiết toàn bộ cách cài đặt, mà là làm rõ vì sao các hướng tiếp cận và công nghệ được lựa chọn phù hợp với bài toán của đồ án. Với mỗi nhóm công nghệ hoặc nền tảng lý thuyết, chương sẽ nêu vấn đề cần giải quyết, các lựa chọn thay thế có thể sử dụng, công nghệ được chọn và lý do lựa chọn. Những chi tiết thiết kế, triển khai và đánh giá hệ thống sẽ được trình bày trong các chương sau.

\section{Mô hình multi-tenant}
\label{section:3.1}

Ở Chương 2, hệ thống được xác định là phục vụ nhiều cửa hàng trên cùng một nền tảng. Mỗi cửa hàng cần có chatbot, nguồn dữ liệu knowledge base, hội thoại, lead và yêu cầu mua hàng riêng. Vì vậy, một nền tảng chỉ phục vụ một cửa hàng đơn lẻ sẽ không phù hợp với phạm vi của đề tài. Mô hình multi-tenant được sử dụng để giải quyết yêu cầu này.

Multi-tenant là mô hình trong đó nhiều tenant cùng sử dụng một hệ thống phần mềm chung, nhưng dữ liệu và cấu hình của từng tenant được tách biệt. Trong phạm vi đồ án, một tenant tương ứng với một cửa hàng nội thất. Mỗi tenant có thể có chatbot, nguồn dữ liệu, hội thoại và các bản ghi nghiệp vụ riêng. Cách tiếp cận này giúp hệ thống không cần triển khai một bản phần mềm độc lập cho từng cửa hàng, đồng thời vẫn đảm bảo mỗi cửa hàng chỉ truy cập dữ liệu thuộc phạm vi của mình.

Có một số hướng tiếp cận thay thế cho bài toán này. Cách thứ nhất là triển khai một hệ thống riêng cho từng cửa hàng. Cách này đơn giản về mặt phân tách dữ liệu, nhưng khó mở rộng và khó bảo trì khi số lượng cửa hàng tăng. Cách thứ hai là dùng chung hoàn toàn dữ liệu giữa các cửa hàng. Cách này đơn giản hơn về thiết kế, nhưng không phù hợp vì dữ liệu sản phẩm, hội thoại và yêu cầu mua hàng của từng cửa hàng cần được tách biệt. Cách thứ ba là dùng mô hình multi-tenant, trong đó hệ thống dùng chung hạ tầng và mã nguồn, nhưng dữ liệu được gắn với tenant tương ứng.

Trong đồ án, mô hình multi-tenant được lựa chọn vì phù hợp với yêu cầu quản lý nhiều cửa hàng trên cùng một nền tảng đã nêu ở Chương 2. Cách tiếp cận này hỗ trợ các chức năng như quản lý tenant, quản lý thành viên cửa hàng, quản lý chatbot theo tenant, quản lý knowledge base theo tenant và phân tách hội thoại, lead, yêu cầu mua hàng giữa các cửa hàng.

\section{Retrieval-Augmented Generation và knowledge base}
\label{section:3.2}

Một yêu cầu quan trọng của hệ thống là chatbot cần trả lời dựa trên dữ liệu sản phẩm hoặc tài liệu riêng của từng cửa hàng, thay vì chỉ dựa vào kiến thức sẵn có của mô hình ngôn ngữ. Nếu chỉ sử dụng mô hình ngôn ngữ lớn mà không bổ sung dữ liệu riêng, chatbot có thể trả lời chung chung hoặc đưa ra thông tin không đúng với dữ liệu của cửa hàng. Vì vậy, đồ án sử dụng hướng tiếp cận Retrieval-Augmented Generation, gọi tắt là RAG.

RAG là hướng tiếp cận kết hợp giữa truy xuất thông tin và sinh ngôn ngữ tự nhiên. Khi người dùng đặt câu hỏi, hệ thống trước tiên truy xuất các đoạn thông tin liên quan từ knowledge base, sau đó đưa các thông tin này vào ngữ cảnh để mô hình ngôn ngữ sinh câu trả lời. Cách tiếp cận này giúp câu trả lời có cơ sở dữ liệu rõ hơn, đồng thời cho phép cập nhật tri thức bằng cách cập nhật knowledge base thay vì huấn luyện lại toàn bộ mô hình. Bài báo của Lewis và cộng sự đã giới thiệu RAG như một cách kết hợp bộ nhớ tham số của mô hình sinh với bộ nhớ phi tham số từ nguồn tri thức bên ngoài~\cite{lewis2020rag}.

Đối với bài toán tư vấn bán hàng nội thất, knowledge base đóng vai trò là nguồn dữ liệu để chatbot truy xuất khi tư vấn. Knowledge base có thể được xây dựng từ thông tin sản phẩm, mô tả sản phẩm hoặc các tài liệu liên quan đến cửa hàng. Trong phạm vi đồ án, knowledge base không chỉ là nơi lưu trữ dữ liệu, mà còn là thành phần giúp đảm bảo chatbot trả lời theo dữ liệu của đúng tenant.

Các lựa chọn thay thế cho RAG gồm: chatbot kịch bản, tìm kiếm từ khóa truyền thống, hoặc fine-tuning mô hình ngôn ngữ. Chatbot kịch bản dễ kiểm soát nhưng khó xử lý hội thoại linh hoạt. Tìm kiếm từ khóa có thể tìm các tài liệu chứa từ tương ứng, nhưng hạn chế khi người dùng diễn đạt nhu cầu bằng ngôn ngữ tự nhiên. Fine-tuning có thể làm mô hình thích nghi với miền dữ liệu cụ thể, nhưng tốn chi phí, khó cập nhật khi dữ liệu sản phẩm thay đổi và không phù hợp với từng cửa hàng nhỏ trong hệ thống. Vì vậy, RAG là lựa chọn phù hợp hơn cho đồ án vì hỗ trợ cập nhật dữ liệu theo từng tenant, phù hợp với hội thoại tự nhiên và giảm phụ thuộc vào việc huấn luyện lại mô hình.

\section{Truy xuất thông tin trong hệ thống RAG}
\label{section:3.3}

Trong hệ thống RAG, truy xuất thông tin là bước quyết định dữ liệu nào được đưa vào ngữ cảnh để chatbot sinh câu trả lời. Đối với dữ liệu sản phẩm nội thất, thông tin có thể bao gồm cả thuộc tính cụ thể như giá, kích thước, chất liệu, loại sản phẩm, và mô tả ngữ nghĩa như phong cách hoặc bối cảnh sử dụng. Vì vậy, hệ thống cần một cơ chế truy xuất phù hợp với cả câu hỏi trực tiếp và câu hỏi diễn đạt tự nhiên.

Có một số hướng truy xuất có thể sử dụng. Hướng thứ nhất là truy xuất theo từ khóa, ví dụ dựa trên từ xuất hiện trong câu hỏi và tài liệu. Hướng này đơn giản, dễ triển khai và phù hợp với các truy vấn chứa tên sản phẩm hoặc thuộc tính rõ ràng. Hạn chế của hướng này là khó xử lý các câu hỏi diễn đạt khác từ nhưng cùng ý nghĩa. Hướng thứ hai là truy xuất theo vector, trong đó văn bản được biểu diễn thành vector ngữ nghĩa và hệ thống tìm các đoạn dữ liệu gần nhất với truy vấn. Hướng này phù hợp hơn với câu hỏi tự nhiên, nhưng cần thêm bước tạo embedding và cơ chế lưu/truy vấn vector. Hướng thứ ba là hybrid retrieval, kết hợp truy xuất từ khóa và truy xuất vector để tận dụng ưu điểm của cả hai cách tiếp cận.

Trong đồ án, truy xuất thông tin được trình bày như một thành phần lý thuyết quan trọng của AI service. Hệ thống có thể sử dụng truy xuất từ khóa, truy xuất vector hoặc kết hợp hai hướng này tùy theo phạm vi triển khai cuối cùng. Điểm quan trọng là kết quả truy xuất phải được giới hạn theo tenant tương ứng, nhằm đảm bảo chatbot của cửa hàng nào chỉ sử dụng dữ liệu của cửa hàng đó khi tư vấn theo cửa hàng.

\section{Mô hình ngôn ngữ lớn trong chatbot tư vấn}
\label{section:3.4}

Mô hình ngôn ngữ lớn được sử dụng để sinh phản hồi tự nhiên cho người dùng sau khi hệ thống đã xác định ngữ cảnh và truy xuất dữ liệu liên quan. Trong đồ án, mô hình ngôn ngữ không hoạt động độc lập, mà là một thành phần trong pipeline RAG. Điều này có nghĩa là câu trả lời của chatbot cần dựa trên dữ liệu đã truy xuất từ knowledge base và lịch sử hội thoại, thay vì chỉ dựa vào kiến thức sẵn có của mô hình.

Đối với bài toán tư vấn nội thất, mô hình ngôn ngữ giúp hệ thống xử lý các yêu cầu bằng tiếng Việt, đặt câu hỏi làm rõ khi thông tin chưa đủ, diễn giải khác biệt giữa các sản phẩm và tạo phản hồi phù hợp với ngữ cảnh hội thoại. Ví dụ, khi người dùng nói rằng muốn tìm một ghế sofa cho phòng khách nhỏ với ngân sách giới hạn, hệ thống cần hiểu các ràng buộc này và phản hồi theo hướng tư vấn thay vì chỉ trả về một danh sách sản phẩm.

Có một số lựa chọn thay thế khi sử dụng mô hình ngôn ngữ. Hệ thống có thể dùng mô hình chạy cục bộ, dùng API từ nhà cung cấp bên ngoài, hoặc kết hợp cả hai tùy điều kiện triển khai. Mô hình cục bộ giúp chủ động hơn về hạ tầng và dữ liệu, nhưng yêu cầu tài nguyên tính toán cao hơn. API mô hình bên ngoài dễ tích hợp và thường có chất lượng tốt, nhưng phụ thuộc vào kết nối mạng, chi phí và giới hạn của nhà cung cấp. Trong phạm vi đồ án, mô hình ngôn ngữ được xem là thành phần có thể thay đổi theo cấu hình, còn trọng tâm của hệ thống là tổ chức pipeline tư vấn, truy xuất knowledge base và quản lý nghiệp vụ multi-tenant.

\section{REST API và kiến trúc giao tiếp giữa các phân hệ}
\label{section:3.5}

Theo yêu cầu ở Chương 2, hệ thống gồm nhiều phân hệ: giao diện web, backend nghiệp vụ, AI service, cơ sở dữ liệu, knowledge base và các kênh chat ngoài. Các phân hệ này cần giao tiếp với nhau theo một cách rõ ràng, dễ kiểm thử và dễ mở rộng. Vì vậy, hệ thống sử dụng hướng giao tiếp qua API, trong đó dữ liệu trao đổi chủ yếu ở định dạng JSON.

REST API là một lựa chọn phổ biến cho các hệ thống web vì tách biệt giao diện người dùng với backend, đồng thời cho phép các chức năng được tổ chức thành các endpoint theo nhóm nghiệp vụ. Trong đồ án, các nhóm API có thể bao gồm đăng nhập, quản lý tenant, quản lý thành viên, quản lý chatbot, quản lý knowledge base, chat, lead, purchase request và thống kê vận hành.

Các lựa chọn thay thế gồm GraphQL, gRPC hoặc giao tiếp trực tiếp giữa frontend và các service con. GraphQL linh hoạt trong truy vấn dữ liệu nhưng có thể phức tạp hơn trong thiết kế và bảo mật đối với phạm vi đồ án. gRPC phù hợp với giao tiếp hiệu năng cao giữa các service, nhưng không thuận tiện bằng REST đối với giao diện web thông thường. Giao tiếp trực tiếp từ frontend đến AI service có thể đơn giản trong demo nhỏ, nhưng không phù hợp vì backend cần kiểm soát xác thực, phân quyền và tenant isolation. Do đó, REST API được lựa chọn vì phù hợp với hệ thống web, dễ triển khai và dễ kiểm thử trong phạm vi đồ án.

\section{Spring Boot cho backend nghiệp vụ}
\label{section:3.6}

Backend nghiệp vụ của hệ thống cần xử lý các chức năng như đăng nhập, phân quyền, quản lý tenant, quản lý thành viên, quản lý chatbot, quản lý nguồn dữ liệu knowledge base, lưu hội thoại, quản lý lead và quản lý yêu cầu mua hàng. Đây là các chức năng có tính nghiệp vụ rõ ràng và cần kết nối với cơ sở dữ liệu quan hệ.

Spring Boot là framework giúp xây dựng các ứng dụng Java độc lập, có thể chạy được với ít cấu hình thủ công và phù hợp với các ứng dụng backend production-grade~\cite{springboot_docs}. Framework này thường được sử dụng cùng hệ sinh thái Spring để xây dựng REST API, quản lý dependency, cấu hình ứng dụng, bảo mật và truy cập cơ sở dữ liệu. Trong đồ án, Spring Boot phù hợp với vai trò backend nghiệp vụ vì hệ thống cần tổ chức các controller, service, repository và entity một cách rõ ràng.

Các lựa chọn thay thế cho backend gồm Node.js với Express/NestJS, Django, Laravel hoặc FastAPI. Node.js phù hợp với ứng dụng JavaScript toàn stack và xử lý I/O tốt. Django có nhiều thành phần tích hợp sẵn cho hệ thống web. FastAPI cũng có thể dùng cho backend, đặc biệt khi toàn bộ hệ thống thiên về Python. Tuy nhiên, trong đồ án này, backend nghiệp vụ và AI service được tách thành hai phần. Spring Boot được lựa chọn cho phần nghiệp vụ vì phù hợp với cấu trúc hệ thống nhiều entity, nhiều role và cần quản lý phân quyền, trong khi FastAPI được dùng cho phần AI service.

\section{FastAPI cho AI service}
\label{section:3.7}

AI service chịu trách nhiệm xử lý các tác vụ liên quan đến hội thoại, truy xuất knowledge base và sinh câu trả lời tư vấn. Thành phần này phù hợp để xây dựng bằng Python vì hệ sinh thái Python có nhiều thư viện hỗ trợ xử lý ngôn ngữ tự nhiên, embedding, retrieval và tích hợp mô hình ngôn ngữ.

FastAPI là framework web hiện đại cho Python, hỗ trợ xây dựng API dựa trên type hints và được thiết kế cho hiệu năng cao~\cite{fastapi_docs}. Trong đồ án, FastAPI phù hợp để xây dựng AI service vì có thể định nghĩa các endpoint nhận yêu cầu từ backend, xử lý truy xuất dữ liệu, gọi mô hình ngôn ngữ và trả kết quả về dưới dạng JSON.

Các lựa chọn thay thế gồm Flask, Django REST Framework hoặc triển khai AI logic trực tiếp trong backend Spring Boot. Flask đơn giản và linh hoạt nhưng cần tự tổ chức nhiều thành phần hơn khi API phức tạp. Django REST Framework mạnh về web backend nhưng không cần thiết nếu AI service chỉ tập trung vào các endpoint xử lý AI. Việc đưa toàn bộ AI logic vào Spring Boot sẽ làm giảm tính tách biệt giữa nghiệp vụ và xử lý AI, đồng thời khó tận dụng hệ sinh thái Python. Vì vậy, tách AI service bằng FastAPI là lựa chọn phù hợp với yêu cầu dễ bảo trì và có khả năng thay đổi mô hình AI trong tương lai.

\section{PostgreSQL cho lưu trữ dữ liệu quan hệ}
\label{section:3.8}

Hệ thống cần lưu nhiều nhóm dữ liệu có quan hệ với nhau, bao gồm tenant, thành viên cửa hàng, chatbot, nguồn dữ liệu knowledge base, hội thoại, tin nhắn, lead, yêu cầu mua hàng và trạng thái vận hành. Các dữ liệu này có cấu trúc tương đối rõ ràng và cần đảm bảo tính nhất quán. Vì vậy, cơ sở dữ liệu quan hệ là lựa chọn phù hợp.

PostgreSQL là hệ quản trị cơ sở dữ liệu quan hệ mã nguồn mở, hỗ trợ SQL và nhiều tính năng phục vụ lưu trữ dữ liệu có cấu trúc~\cite{postgres_docs}. Trong đồ án, PostgreSQL được sử dụng để lưu dữ liệu nghiệp vụ chính của hệ thống. Việc sử dụng PostgreSQL giúp hệ thống biểu diễn rõ các quan hệ như tenant có nhiều chatbot, chatbot có nhiều hội thoại, hội thoại có nhiều tin nhắn, hội thoại có thể phát sinh lead hoặc yêu cầu mua hàng.

Các lựa chọn thay thế gồm MySQL, MongoDB hoặc SQLite. MySQL cũng là một hệ quản trị cơ sở dữ liệu quan hệ phổ biến và có thể đáp ứng nhiều yêu cầu tương tự. MongoDB linh hoạt với dữ liệu dạng tài liệu, nhưng với hệ thống có nhiều quan hệ nghiệp vụ và yêu cầu phân quyền theo tenant, cơ sở dữ liệu quan hệ dễ kiểm soát hơn. SQLite phù hợp với demo nhỏ hoặc ứng dụng cục bộ, nhưng không phù hợp bằng PostgreSQL khi hệ thống có nhiều service và nhiều người dùng. Vì vậy, PostgreSQL được lựa chọn để lưu trữ dữ liệu nghiệp vụ chính.

\section{React cho giao diện người dùng}
\label{section:3.9}

Giao diện người dùng của hệ thống gồm các nhóm màn hình như đăng nhập, quản trị hệ thống, quản trị cửa hàng, quản lý chatbot, quản lý knowledge base, giao diện chat, danh sách hội thoại, lead và yêu cầu mua hàng. Các màn hình này cần tương tác với backend qua API và cập nhật trạng thái giao diện theo hành động của người dùng.

React là thư viện JavaScript dùng để xây dựng giao diện người dùng từ các component độc lập~\cite{react_docs}. Cách tiếp cận component phù hợp với hệ thống vì nhiều phần giao diện có thể được tái sử dụng, ví dụ bảng danh sách, form tạo/cập nhật, khung chat, trạng thái tải và thông báo lỗi. Với giao diện chat, React cũng phù hợp vì cần cập nhật danh sách tin nhắn khi người dùng gửi câu hỏi hoặc khi hệ thống trả lời.

Các lựa chọn thay thế gồm Vue, Angular hoặc xây dựng giao diện bằng template server-side. Vue có cú pháp đơn giản và dễ học, Angular có cấu trúc đầy đủ cho ứng dụng lớn, còn server-side template phù hợp với các hệ thống ít tương tác. Trong phạm vi đồ án, React là lựa chọn phù hợp vì hỗ trợ xây dựng giao diện web tương tác, tách biệt với backend và dễ tổ chức theo component.

\section{Xác thực và phân quyền bằng token}
\label{section:3.10}

Hệ thống có nhiều vai trò khác nhau như quản trị hệ thống, quản trị cửa hàng và nhân viên cửa hàng. Các vai trò này có quyền truy cập khác nhau, ví dụ quản trị hệ thống quản lý tenant, quản trị cửa hàng quản lý chatbot và knowledge base của tenant, còn nhân viên cửa hàng xử lý lead và yêu cầu mua hàng. Vì vậy, hệ thống cần cơ chế xác thực và phân quyền rõ ràng.

Một hướng phổ biến trong các hệ thống web là sử dụng token để xác thực người dùng sau khi đăng nhập. JSON Web Token, gọi tắt là JWT, là một chuẩn biểu diễn claims giữa các bên dưới dạng compact và URL-safe; các claims có thể được ký hoặc bảo vệ toàn vẹn~\cite{jwt_rfc}. Trong hệ thống, token có thể được dùng để backend xác định người dùng, vai trò và phạm vi tenant khi xử lý yêu cầu.

Các lựa chọn thay thế gồm session lưu phía server, OAuth2 đầy đủ hoặc API key. Session phía server dễ hiểu và phù hợp với ứng dụng web truyền thống, nhưng cần quản lý trạng thái phiên trên server. OAuth2 phù hợp với hệ thống tích hợp đăng nhập bên thứ ba hoặc phân quyền phức tạp, nhưng có thể vượt quá phạm vi đồ án. API key phù hợp với giao tiếp giữa hệ thống và dịch vụ, nhưng không phù hợp để đại diện cho từng người dùng nội bộ. Vì vậy, token-based authentication là lựa chọn phù hợp cho hệ thống web có frontend tách biệt backend và cần phân quyền theo vai trò.

\section{Tích hợp Messenger và Telegram bằng webhook}
\label{section:3.11}

Ngoài giao diện web chat, hệ thống còn hướng tới việc tích hợp các kênh chat ngoài như Messenger và Telegram. Mục tiêu của việc tích hợp là cho phép người dùng cuối trao đổi với chatbot qua các kênh quen thuộc, trong khi hệ thống vẫn xử lý nội dung bằng cùng pipeline tư vấn sản phẩm theo cửa hàng.

Webhook là cơ chế trong đó nền tảng bên ngoài gửi yêu cầu HTTP đến hệ thống khi có sự kiện, ví dụ khi người dùng gửi tin nhắn. Telegram Bot API là một giao diện HTTP cho phép xây dựng bot Telegram và hỗ trợ việc nhận cập nhật từ người dùng~\cite{telegram_bot_api}. Với Messenger, hệ thống cũng có thể triển khai theo hướng nhận sự kiện tin nhắn qua webhook, xác định chatbot hoặc tenant tương ứng, xử lý nội dung và gửi phản hồi lại kênh ngoài.

Các lựa chọn thay thế cho webhook là polling hoặc chỉ hỗ trợ web chat nội bộ. Polling đơn giản trong một số trường hợp nhưng kém hiệu quả hơn vì hệ thống phải liên tục kiểm tra cập nhật mới. Chỉ hỗ trợ web chat sẽ giảm độ phức tạp, nhưng chưa đáp ứng yêu cầu tích hợp các kênh mà cửa hàng có thể sử dụng trong thực tế. Vì vậy, webhook là hướng phù hợp cho tích hợp Messenger và Telegram trong phạm vi đồ án.

\section{Docker cho đóng gói và triển khai}
\label{section:3.12}

Hệ thống gồm nhiều thành phần như frontend, backend, AI service và PostgreSQL. Nếu chạy từng thành phần trực tiếp trên máy, việc cài đặt môi trường có thể phức tạp và dễ phát sinh lỗi do khác biệt phiên bản thư viện, hệ điều hành hoặc cấu hình. Vì vậy, đồ án sử dụng hướng đóng gói bằng container để hỗ trợ triển khai và chạy thử.

Docker cho phép đóng gói ứng dụng cùng các phụ thuộc cần thiết vào container, giúp ứng dụng chạy trong môi trường tách biệt và nhất quán hơn giữa các máy~\cite{docker_docs}. Trong phạm vi đồ án, Docker phù hợp để chạy các service như backend, AI service và cơ sở dữ liệu trong môi trường demo. Khi kết hợp với Docker Compose, các service có thể được khai báo và khởi chạy cùng nhau, giúp quá trình demo và kiểm thử thuận tiện hơn.

Các lựa chọn thay thế gồm cài đặt thủ công toàn bộ môi trường hoặc triển khai trực tiếp lên máy chủ không dùng container. Cài đặt thủ công có thể phù hợp khi hệ thống rất nhỏ, nhưng khó tái lập môi trường. Triển khai trực tiếp không dùng container có thể tối ưu hơn trong một số môi trường production, nhưng yêu cầu quản trị hệ thống nhiều hơn. Với đồ án tốt nghiệp, Docker là lựa chọn phù hợp vì giúp đóng gói, chạy thử và tái lập môi trường dễ dàng.

\section{Kết chương}
\label{section:3.ket_chuong}

Chương này đã trình bày các nền tảng lý thuyết và công nghệ chính được sử dụng trong đồ án. Mô hình multi-tenant được lựa chọn để đáp ứng yêu cầu phục vụ nhiều cửa hàng trên cùng một nền tảng, trong khi vẫn tách biệt dữ liệu và cấu hình của từng tenant. RAG và knowledge base được sử dụng để giúp chatbot trả lời dựa trên dữ liệu riêng của cửa hàng, phù hợp với bài toán tư vấn sản phẩm nội thất. Các nội dung về truy xuất thông tin và mô hình ngôn ngữ lớn làm rõ cách hệ thống xử lý câu hỏi tự nhiên và tạo câu trả lời tư vấn theo ngữ cảnh.

Bên cạnh đó, chương đã trình bày các công nghệ phục vụ triển khai hệ thống, bao gồm Spring Boot cho backend nghiệp vụ, FastAPI cho AI service, PostgreSQL cho lưu trữ dữ liệu quan hệ, React cho giao diện người dùng, token-based authentication cho xác thực và phân quyền, webhook cho tích hợp Messenger/Telegram, và Docker cho đóng gói triển khai. Với mỗi công nghệ, chương đã nêu vấn đề cần giải quyết, một số lựa chọn thay thế và lý do lựa chọn trong phạm vi đồ án. Các nền tảng và công nghệ này là cơ sở để trình bày thiết kế, triển khai và đánh giá hệ thống trong các chương tiếp theo.

\end{document}